\documentclass[10pt,a4paper]{amsart}
\usepackage[margin=19mm]{geometry}
\usepackage{amsmath,amssymb,mathtools}
\usepackage[scr=rsfs,cal=euler]{mathalfa}
\usepackage{xfrac}
\usepackage[unicode,hidelinks,bookmarks=false]{hyperref}
\newcommand{\ie}{i.e.\ }
\newcommand{\cf}{cf.\ }
\newcommand{\resp}{respectively\ }
\newcommand{\Subsection}{\subsection}
\providecommand{\nobreakdash}{\nobreak\hbox{-}}
\setlength{\emergencystretch}{2em}
\multlinegap0pt
\newtheorem{theorem}{Theorem}
\newcommand{\ByteString}{\mathsf{ByteString}}
\newcommand{\ByteStringBounded}{\ByteString_{w}}
\newcommand{\BytesVar}{\mathit{bytes}}
\newcommand{\DecodeM}{\mathsf{Decode}_m}
\newcommand{\EncodeM}{\mathsf{Encode}_m}
\newcommand{\Fin}{\mathsf{Fin}}
\newcommand{\LeanPayload}[2]{\href{https://github.com/KyrylR/phd-symmetric-cryptography/blob/8e903222a7df1c4365592e4b16fa752a924b8977/lean/encoding-dr-1/EncodingProofs/BaseMLen/Payload.lean\#L#1-L#2}{[Payload.lean, L#1--L#2]}}
\newcommand{\LeanStream}[2]{\href{https://github.com/KyrylR/phd-symmetric-cryptography/blob/8e903222a7df1c4365592e4b16fa752a924b8977/lean/encoding-dr-1/EncodingProofs/BaseMLen/Stream.lean\#L#1-L#2}{[Stream.lean, L#1--L#2]}}
\newcommand{\List}{\mathsf{List}}
\newcommand{\OK}{\mathsf{ok}}
\newcommand{\concat}{\mathbin{\|}}
\providecommand{\href}[2]{#2}
\title{Canonical Byte-String Encoding for Finite-Ring Cryptosystems}
\author{Kyrylo Riabov, Serhii Kryvyi}
\date{Source: arXiv:2603.23364v2 (CC BY 4.0)}
\begin{document}
\maketitle

\noindent\emph{Source and licence.} Canonical Byte-String Encoding for Finite-Ring Cryptosystems. Version arXiv:2603.23364v2 (CC BY 4.0), distributed by arXiv under the Creative Commons Attribution 4.0 International licence, http://creativecommons.org/licenses/by/4.0/, as recorded in the arXiv licence metadata for this exact version and shown on the abstract page https://arxiv.org/abs/2603.23364v2. Full licence notice: \texttt{licenses/arxiv-2603.23364-CC-BY-4.0.txt}.\par\smallskip

\noindent\textit{Editorial scope.} Counted unit: the article's theorem thm:stream-roundtrip (source0.tex lines 440--451). The article's complete original proof of it is delivered with it (source0.tex lines 453--463, one \texttt{proof} environment of 11 lines). No mathematics is written, repaired or re-derived: the statement and the proof are the article's own lines, in the article's own notation, and no retained line is substituted or re-flowed.  No further source-local context is needed: every cross-reference of every retained line resolves inside the two delivered spans.  Nothing was omitted from any retained span, so the package carries no omission record.  No retained line contains a citation, so the wrapper carries no bibliography and no bibliography entry.  No retained line contains a figure environment, an image-inclusion command or drawing code, so no image is delivered and the package declares no asset dependency.  Editorial formatting only, in addition: an AMS \texttt{amsart} preamble that restates the article's own declarations and macros used by the retained lines, namely the environments \texttt{theorem} on separate counters, together with the standard packages those lines need.  The retained environments print in this document as Theorem 1 in order of appearance, whereas the article prints the same items with the numbering of its own full document.  The complete native source of the article as distributed in the arXiv e-print for this exact version is bundled unchanged as \texttt{sources/197078/source0.tex}.
    \begin{theorem}[Stream correctness]
        \label{thm:stream-roundtrip}
        For every supported modulus \(m\) and every admissible byte string
        \(\BytesVar \in \ByteStringBounded\),
        \[
            \DecodeM(\EncodeM(\BytesVar)) = \OK(\BytesVar).
        \]
        Moreover, for every valid suffix \(t \in \List(\Fin(m))\),
        \[
            \DecodeM(\EncodeM(\BytesVar) \concat t) = \OK(\BytesVar).
        \]
    \end{theorem}

    \begin{proof}
        In the Lean 4 formalization, the suffix-tolerant theorem is \LeanStream{174}{254}.
        Its empty-suffix specialization is \LeanStream{256}{266}.
        The proof uses the header-parsing lemmas in \LeanStream{35}{133}.
        The zero-length lemmas appear in \LeanStream{135}{172}, and the payload roundtrip theorem in
        \LeanPayload{79}{141}.
        It first recovers the declared length prefix, then the stored state prefix, and finally applies the
        payload roundtrip theorem to the remaining suffix.
        The extra-suffix clause follows because decoding stops after reconstructing the declared number
        of bytes and therefore leaves any remaining valid suffix unread.
    \end{proof}

\end{document}
